\section{Future Research}

Five directions emerge naturally from the scope boundaries of this study:

\begin{enumerate}
  \item \textbf{Autoscaling response under real spike workloads} --- extend the study to a
        multi-node cluster with \textsc{hpa} enabled and measure the scale-out window and
        its job-level impact \citep{tran2022, yuan2024}.
  \item \textbf{Replication across orchestration frameworks} --- test whether the crossover
        point and degradation curve profile hold for Airflow and Prefect under equivalent
        workload conditions.
  \item \textbf{Cloud-hosted infrastructure} --- repeat the experiments on \textsc{gcp}
        \textsc{gke} or \textsc{aws} \textsc{eks} to determine whether cloud network
        overhead shifts the crossover point.
  \item \textbf{Real workloads with variable resource demand} --- replace synthetic matrix
        computations with real \textsc{etl} or data processing pipelines where resource
        consumption is unpredictable.
  \item \textbf{Multi-tenant environments} --- investigate how the crossover point changes
        when multiple teams share the same Kubernetes cluster and node resources are
        contended at the cluster level rather than the pod level.
\end{enumerate}
